\chapter*{Leitidee}
\addcontentsline{toc}{chapter}{Leitidee}

\textbf{Mathematik 1} entwickelt die formale Sprache und die grundlegenden Strukturen f�r die weitere mathematische Analyse.

Wir beginnen mit Mengen und Logik zur pr�zisen Beschreibung mathematischer Aussagen. �ber Zahlenbereiche, Vektorr�ume und lineare Abbildungen gelangen wir zu Funktionen als zentralem Objekt der Analysis. Funktionen beschreiben Abh�ngigkeiten zwischen Gr��en und bilden damit eine wesentliche Grundlage f�r die mathematische Modellierung technischer Zusammenh�nge. Im Mittelpunkt stehen zun�chst lokale Eigenschaften von Funktionen, die wir mit und ohne Differentialrechnung untersuchen.

Ziel dieses ersten Teils ist es, ein strukturelles Fundament zu schaffen: mathematische Objekte klar zu definieren, Beziehungen pr�zise zu formulieren und Funktionen als Werkzeug zur Beschreibung und Analyse technischer Zusammenh�nge zu verstehen.

%\chapter*{Leitidee}
%\addcontentsline{toc}{chapter}{Leitidee}
%
%\textbf{Mathematik 1} entwickelt die formale Sprache und die grundlegenden Strukturen, auf denen die weitere mathematische Analyse aufbaut.
%
%Wir beginnen mit Mengen und Logik als pr�ziser Beschreibung mathematischer Aussagen. Darauf aufbauend f�hren wir �ber Zahlenbereiche, Vektorr�ume und lineare Abbildungen zu Funktionen. Funktionen sind das zentrale Objekt der Analysis und wir lernen vor allem lokale Eigenschaften von Funktionen mit und ohne Differentialrechnung zu untersuchen.
%
%Ziel dieses ersten Teils ist es, ein \textbf{strukturelles Fundament} zu schaffen: mathematische Objekte klar zu definieren, Beziehungen pr�zise zu formulieren und Funktionen in ihrer lokalen Struktur zu verstehen.
